\section{Principes variationnels}

    \subsection{Principe de Fermat}

    Rayon optique choisit la trajectoire qui lui prend le moins de temps : « principe d’économie naturelle ». 

    \subsection{Principe de Maupertuis}



    \subsection{Principe de moindre action}

    Lagrange proposa un changement de point de vue sur la mécanique : plutôt que de déterminer la position $\mathbf{r}(t)$ et la vitesse $\mathbf{v}(t)$ d’une particule connaissant son état initial $(\mathbf{r}(0), \mathbf{v}(0))$, on s’intéresse à la trajectoire effectivement suivie par une particule partant de $\mathbf{r}_1$ à $t_1$ et arrivant en $\mathbf{r}_2$ en $t_2$. 

    \begin{omed}{Principe 1 (lagrangien)}
        Tout système mécanique est caractérisé par un \textbf{lagrangien} $\mathcal{L}(x, \dot{x}, t)$ où $x$ et $\dot{x}$ sont appelées \textbf{variables d’état}.
    \end{omed}

    \begin{omed}{Principe 2 (dit de moindre action)}
        Pour toute trajectoire $x(t)$ entre $(x_1, t_1)$ et $(x_2, t_2)$, on définit l’\textbf{action} $S$ par 
        \[ S = \int_{t_1}^{t_2} \mathcal{L}(x,\dot{x},t)dt \]  
        Le principe de moindre action dit que la trajectoire physique effectivement suivie $X(t)$ est telle que $S$ est minimale.
    \end{omed}

